\section{经济与社会：霸权背后的土地、劳役与交通}

春秋的会盟和出兵，常在《春秋》中只留下一个“会”“伐”或“城”字。字很短，成本很长：粮食要离开田亩，车马要过河渡和山口，守城者与筑城者又不能同时耕作。传世材料没有赋税簿、户籍册和粮仓清册，因而不能给任何诸侯国算出收入或兵额；不过，战争、救荒、迁城和工程反复进入编年，足以让人看见政治秩序怎样不断向土地、人力和道路索取。

\subsection{灾年进入年表，不能自动变成经济统计}

自然异象之所以值得放在经济社会史中，并不是它们能替我们量出产量，而是编年者反复选择了“记下”它们。前714年，鲁历先记大雨震电，继而记大雨雪，\eventanchor{SA-0714-RAIN-SNOW}{春秋·鲁历大雨震电与大雨雪}；前657年冬春不雨，\eventanchor{SA-0657-DROUGHT}{春秋·鲁历冬春不雨}正处于齐、宋会盟及南向筹备的年代；前625年又有自冬至秋七月不雨的长时段记录，\eventanchor{SA-0625-DROUGHT}{春秋·鲁历久旱}。这些经文能够证明鲁国编年传统把异常天气视为应当登记的事，不能告诉我们旱区的范围、谷物的减产比例，或哪一位诸侯因此改变了一项既定战略。

同样的限度适用于前539年大雨雹，\eventanchor{SA-0539-HAIL}{春秋·鲁历大雨雹}、前523年地震，\eventanchor{SA-0523-EARTHQUAKE}{春秋·鲁历地震}、前509年陨霜杀菽，\eventanchor{SA-0509-FROST-BEANS}{春秋·鲁历陨霜杀菽}以及前492年地震，\eventanchor{SA-0492-EARTHQUAKE}{春秋·鲁历地震}。它们都是鲁国年表中的可见痕迹，不是覆盖列国的气象台记录。现代古气候、环境史和地震史研究可以帮助追问较宽时段的环境条件，却不能让某一条“雨雹”或“地震”直接证明全国饥荒，更不能把灾异按后世的报应逻辑写成政变或战败的原因。

不过，编年中的饥、螽与火仍提醒我们，粮食、居住空间和防灾并非政治叙事外的背景。前564年，\eventanchor{SA-0564-SONG-FIRE}{春秋·宋国火灾}；前557年鲁历又记\eventanchor{SA-0557-EARTHQUAKE}{春秋·鲁历地震}；前533年，\eventanchor{SA-0533-CHEN-FIRE}{春秋·陈国火灾}；前525年宋、卫、陈、郑同年见灾，\eventanchor{SA-0525-CENTRAL-PLAINS-FIRES}{春秋·中原诸国同年见灾}；前483年鲁历记螽灾，\eventanchor{SA-0483-LOCUSTS}{春秋·鲁历螽灾}，前481年又记饥，\eventanchor{SA-0481-FAMINE}{春秋·鲁历饥}。这些记载不足以说明灾害怎样沿道路传播、各国是否同受影响，甚至“同年见灾”本身也不能推出同一种自然原因；但它们让我们看见，诸侯在盟会、征战和城防之外，还须面对农时失调、物资短缺和地方承受力的持续压力。

\subsection{田亩进入征调}

前594年鲁国\eventanchor{SA-0594-INITIAL-LAND-TAX}{春秋·鲁初税亩}。《春秋》的年次和《左传》《国语》关于初税亩的传说，至少使“田亩”成为观察负担的一处可见入口：执政者试图把征取同可辨认的受田联系起来。它并不证明此前没有任何征收，也不能由一条传文推出整齐的井田崩解、固定的什一税率或全国统一的地主结构。更可把握的是，军事和政务的需求，使土地的占有和承担开始被更明确地放在同一张账上。

六年后，\eventanchor{SA-0590-LU-QIUJIA}{春秋·鲁作丘甲}，又有\eventanchor{SA-0590-LU-DEFENSE-REPAIR}{春秋·鲁修赋缮完具守备}的叙事。丘、甲究竟是怎样的土地、户口或兵役单位，后世解释并不一致，不能折算为准确兵员；但在齐楚结好、鲁须自保的处境中，田地、甲士和城防确被连了起来。到前541年，\eventref{SA-0541-ZICHAN-LAND}{春秋·子产整顿田制赋役与城防}所保存的是小国郑试图让征调和守备较为可预期的记忆，而不是一部可复原条文的财政法典。前483年鲁\eventanchor{SA-0483-LU-TIANFU}{春秋·鲁用田赋}，再次以田亩组织军赋；它没有挽回公室在三桓面前的弱势，反而提示征调的难题已同谁掌握邑兵、家臣和地方资源纠缠在一起。

因此，土地不能只写成农业背景。鲁的初税亩、丘甲和田赋是几个有年次的制度切口；晋卿、齐田氏和鲁三桓的资源积累则主要由经传叙事、盟书和后出的国别史拼合而成。春秋末侯马盟书可旁证晋系精英以书面盟誓处理结纳与排斥，却不能把某一处盟书的格式推广为各国都已有同样的行政体系。国君借卿族征兵、卿族借军政职位积累地方关系，这一相互依赖比“贵族取代国家”的整齐公式更接近材料所能支持的图景。

\subsection{点兵、裁军与大蒐：谁能把人带到一处}

《春秋》保存的“阅”“作军”“舍军”和“大蒐”，让我们偶尔能看见动员从抽象的“有兵”变成聚集在一处的人和器物。前706年鲁于八月大阅，\eventanchor{SA-0706-LU-MUSTER}{春秋·鲁国大阅}，能够确定的是一次检阅军队的行动进入了鲁历；队伍人数、兵种比例、训练内容和普通家庭的服役期限都没有留下。前562年鲁作三军，\eventanchor{SA-0562-LU-THREE-ARMIES}{春秋·鲁作三军}，前537年又舍中军，\eventanchor{SA-0537-LU-ABOLISH-CENTRAL-ARMY}{春秋·鲁舍中军}。两事相隔多年，足以说明军政分配会被调整，不能被简单折算成一条从“扩军”到“裁军”的现代国防预算曲线。

前531年鲁大蒐于比蒲，\eventanchor{SA-0531-LU-BIPU-HUNT}{春秋·鲁大蒐于比蒲}；前497年又筑蛇渊囿并大蒐于比蒲，\eventanchor{SA-0497-SNAKE-LAKE}{春秋·鲁筑蛇渊囿与大蒐}。大蒐兼有礼仪、田猎和军事检阅的面向，传世叙事虽能让人看见贵族、车马和随从在场，却不能据此计算常备军规模。蛇渊囿的营建也不应被写成一项孤立的君主奢侈：它要求劳力、地面和组织，所见的是城邑周边空间能被政治权力占用；谁实际承担了多少劳役，现有材料仍无法开列。

这些零散节点和初税亩、丘甲互相照亮。它们并不证明鲁已经建立一套以户籍直接统率的军事机器，却说明田地、城邑、礼仪与可被召集的人手愈发难以分开。国家若要应付外部威胁，须借卿族和地方网络调动资源；这些网络也正因承担动员而积累了独立的分量。

\subsection{城墙、道路与交换}

城并非只是在地图上添一个据点。前571年诸侯会后\eventanchor{SA-0571-HULAO}{春秋·诸侯城虎牢}，晋以郑东部要道上的城防施压，齐、宋、卫、曹以及邾、滕、薛等参与者也被卷入工程。经文能确认筑城的年次，叙事能显示其战略方向；谁交出多少木石、粮食和劳力，却没有留存。前510年诸侯大夫又\eventanchor{SA-0510-CHENGZHOU-WALL}{春秋·诸侯城成周}。周室内战后的都城防务需要诸侯共同负担，这一事实比“王室尚有无上权威”的抽象判断更具体：王室仍可动员名分，工程资源却不再能够独自供给。

城址、聚落和道路的考古调查能够补足传世编年的空间感，显示城邑层级、通行条件和区域环境；它们不能替我们填出虎牢或成周的工期、墙长和徭役人数。墓葬、车马坑、兵器与铸造遗存同样说明车战、水陆行动和精英消费有物质条件，却不是普通家庭收入或市场规模的抽样。

所谓“贸易”尤其须收紧笔墨。现有材料较清楚地记录的是粮食、礼物、车马和军队沿交通线的移动，而非连续的价格、税关或商人名册。前647年\eventref{SA-0647-QIN-JIN-GRAIN}{春秋·秦输粟援晋}是一次饥荒中的跨国输送：它既关乎灾年的生计，也是一段盟约关系的信用，不应轻易改写成自由市场的谷物流通。秦军东向郑而在殽道受挫的叙事，则提醒关中通向中原的河渡、谷道既可交换也可截断；关于郑商人的故事和每一车粮的路线，均主要依赖后出的叙事，不能当作当时的运输档案。春秋城市当然是物资和人群交会之地，但我们尚不能据此断言存在一个由国家常设管理、覆盖列国的统一市场。

\subsection{修一座城，也是在修复谁能进入它}

城防工程的成本不只在边境。前508年鲁雉门及两观失火，\eventanchor{SA-0508-PALACE-GATES-FIRE}{春秋·鲁雉门两观失火}，同年又新作雉门及两观，\eventanchor{SA-0508-REBUILD-PALACE-GATES}{春秋·鲁重建雉门两观}。经文与《左传》让我们知道失火与重建相继发生，却不保存火因、建筑尺度、木料来源或修筑者名单。两观和宫门所关涉的是公室可见的出入与礼仪空间；重建既是物资和劳力的调度，也是把受损的公室秩序重新摆到城中。

前504年鲁城中城，\eventanchor{SA-0504-LU-WALL-ZHONGCHENG}{春秋·鲁城中城}；前495年又城漆，\eventanchor{SA-0495-LU-WALL-QI}{春秋·鲁城漆}；前492年季孙斯、叔孙州仇城启阳，\eventanchor{SA-0492-LU-WALL-QIYANG}{春秋·季孙叔孙城启阳}。这些短记不能告诉我们城墙是土是石、劳役如何分派，或工程完成后谁拥有税赋和军权；它们却使“筑城”不能只被当作地图上的一笔。不同的城、不同的主持者，意味着公室、卿族与地方据点都须把劳力、物资和守备组织到一处。尤其由三桓主导的启阳工程，不能直接证明鲁国所有城防已经私家化，却提示工程所需的资源并不必然只通过国君一条渠道汇集。

\subsection{迁徙不是空白处的箭头}

战争把人从土地上移开，也把新的劳役和守备问题带到迁入地。前659年\eventref{SA-0659-QI-XING}{春秋·齐迁邢于夷仪}，传世经年保存了狄人进攻与救援的骨架，《左传》才把迁徙置于夷仪。旧城、田地和宗庙不能原样搬运，齐的援助也把筑城、戍守和未来会盟的期待联系起来；人数、安置地块和迁城工程的细目却无从复原。前660年卫国旧都受创后，\eventref{SA-0660-WEI-DAI}{春秋·宋接卫遗民}所说的渡河、立君与齐军戍守，同样只容许我们确认一个国家名号、部分人口和防卫网络被重新接续，不能画成完整的人口流向图。

楚、吴和诸侯在江淮与中原边缘的灭国、迁许与迁蔡，显示迁徙并非北方独有的现象。可是《春秋》多记“迁”而不记迁民，城址的年代也常只给出宽阔区间；迁都由谁主导、旧城怎样处置、原有人群是否留在当地，都必须逐案讨论。把迁徙写成一条从战乱自动通向“开发”的道路，恰好会抹掉材料最重要的沉默。

\subsection{不同腹地，不同成本}

齐向海岱东部扩展并\eventanchor{SA-0567-QI-LAI-FALL}{春秋·齐灭莱}，得到的是更深的东部城邑、人口和盐海资源的可能性；经传记录灭国，考古和区域研究才能追问整合的空间条件，却仍不能说明莱民如何被安置。秦的关中腹地面对的是河渡、山口和东进距离，输粟、围郑与殽道的事件链使其后勤约束格外清楚。中原小国夹在晋楚道路之间，城防、会盟和临时征调更容易直接落到有限的邑地；郑的田制整顿正应放在这样的尺度上看。

至于楚的江汉、吴越的江淮水网，城址、兵器和带铭金属器能显示水陆交通、铸造能力与精英网络并不等同于中原，却不能由器物推出航运吨数、商业税或某次战役的水路部署。区域差异不是一张“富强诸侯”的排行榜，而是不同地形、城邑和政治网络令征粮、筑城、迁民与出兵付出不同代价。

\begin{sourcenote}[文本的账与考古的尺度]
《春秋》最适合固定鲁国编年中可见的灾异、筑城、会盟和迁徙；《左传》《国语》提供征发、救援和田制的叙事关联，但不是现场的工程簿或谈判记录；《史记》又以汉代的国别兴亡线重组早期传统。城址、盟书、墓葬和兵器能够校正空间、技术和精英网络的想象，却不能补出失传的人口、税额和市场数据。两类材料互相提问，不能相互代写。
\end{sourcenote}

\begin{turningpoint}[从春秋社会到战国国家]
\term{它解决了什么}：以封邑、宗族、盟约和临时征调组织资源，能让诸侯在没有完备常设官僚的条件下筑城、救荒和出兵。

\term{谁承担代价}：土地占有者、服役者、迁城的人群和小国盟友承担最直接的粮食、人力与安全成本；强卿与受援国也可借此积累自己的关系和据点。

\term{它留下什么压力}：资源多经由家族和邑地流动，国君未必能把赋役直接收归公室。前483年\eventref{SA-0483-LU-TIANFU}{春秋·鲁用田赋}所显示的重新联结田亩与军赋，正预示战国诸国会把土地、户籍、军役和税收更直接地接到国家手中；这不是春秋已经完成的制度，而是春秋动员困境留下的问题。
\end{turningpoint}
